
c'est un anneau d'intersection complète (EGA IV 19.3.2), et par suite @@M0@@ est d'intersection complète en @@M1@@ au sens de EGA IV 19.3.6.

Réciproquement, supposons que @@M2@@ soit un morphisme d'intersection complète au sens de EGA IV 19.3.6, et soit @@M3@@. Soit @@M4@@ un voisinage affine de @@M5@@, au-dessus d'un voisinage affine @@M6@@ de @@M7@@. On peut alors trouver un @@M8@@-schéma lisse @@M9@@, et une immersion fermée @@M10@@ (en prenant par exemple pour @@M11@@ un fibré vectoriel sur @@M12@@). Nous allons montrer que @@M13@@ est régulière. Or d'après EGA IV 19.3.7, @@M14@@ est transversalement régulier en @@M15@@, puisque @@M16@@ est lisse sur @@M17@@, donc @@M18@@ régulier en @@M19@@, et @@M20@@ plat sur @@M21@@. Il résulte alors de EGA IV 19.2.4, a)@@M22@@c), que @@M23@@ est régulière en @@M24@@. Par suite, @@M25@@ est un morphisme d'intersection complète au sens de 1.1.

\textbf{Proposition 1.5.} \uline{Soient} @@M26@@, \uline{et} @@M27@@ \uline{deux morphismes localement d'intersection complète}. \uline{Alors} @@M28@@ \uline{est un morphisme localement d'intersection complète}.

Soit @@M29@@. On peut trouver des voisinages ouverts @@M30@@ de @@M31@@, et des immersions régulières @@M32@@ telles qu'on ait les factorisations :
@@M33@@
On peut compléter ce diagramme par le carré commutatif :
